% =====================================================================
%  HR-ETE-GNN: Robust Volatility Forecasting via Hurst-Regime-Adaptive
%  Rényi Effective Transfer Entropy Graph Neural Networks
%
%  Undergraduate thesis -- Lorenzo Uriel L. Ramoy
%  Converted from [RAMOY]Thesis_Manuscript.docx
%
%  Compile with pdfLaTeX (Overleaf default).
% =====================================================================
\documentclass[12pt,oneside]{report}

% ---------- Encoding & fonts ----------
\usepackage[T1]{fontenc}
\usepackage[utf8]{inputenc}
\usepackage{amsmath,amssymb}       % must load before newtxmath
\usepackage{newtxtext,newtxmath}   % Times New Roman look, as in the .docx
\usepackage{microtype}
\setlength{\emergencystretch}{3em}

% ---------- Page layout ----------
\usepackage[letterpaper,margin=1in]{geometry}
\usepackage{setspace}
\doublespacing                     % the .docx uses double spacing

% ---------- Math ----------
\DeclareMathOperator{\QLIKE}{QLIKE}
\DeclareMathOperator{\ETE}{ETE}
\DeclareMathOperator{\ERTE}{ERTE}
\DeclareMathOperator{\ReLU}{ReLU}

% ---------- Tables, lists ----------
\usepackage{booktabs}
\usepackage{array}
\usepackage{enumitem}
\usepackage{caption}
\captionsetup[table]{labelfont=bf,textfont=bf,justification=raggedright,
                     singlelinecheck=false,labelsep=space}

% ---------- Headings: "CHAPTER I" / "INTRODUCTION" ----------
\usepackage{titlesec}
\renewcommand{\thechapter}{\Roman{chapter}}
\renewcommand{\thesection}{\arabic{chapter}.\arabic{section}}
\titleformat{\chapter}[display]
  {\normalfont\bfseries\centering}
  {\MakeUppercase{\chaptertitlename}\ \thechapter}
  {0pt}
  {\MakeUppercase}
\titlespacing*{\chapter}{0pt}{0pt}{24pt}
\titleformat{\section}{\normalfont\bfseries}{\thesection}{0.5em}{}
\titleformat{\subsection}{\normalfont\bfseries}{\thesubsection}{0.5em}{}
\titleformat{\subsubsection}{\normalfont\bfseries}{\thesubsubsection}{0.5em}{}
\setcounter{secnumdepth}{2}
\setcounter{tocdepth}{2}

% Table numbers follow the section, e.g. "Table 3.5.1"
\usepackage{chngcntr}
\counterwithin{table}{section}
\counterwithin{figure}{section}

% ---------- Front-matter headings ----------
\renewcommand{\contentsname}{TABLE OF CONTENTS}
\renewcommand{\listfigurename}{LIST OF FIGURES}
\renewcommand{\listtablename}{LIST OF TABLES}
\renewcommand{\bibname}{BIBLIOGRAPHY}
\renewcommand{\appendixname}{APPENDIX}

% ---------- Links ----------
\usepackage[hidelinks]{hyperref}
\usepackage{bookmark}

% Unnumbered chapter-level heading that still appears in the ToC
\newcommand{\frontchapter}[1]{%
  \chapter*{#1}\addcontentsline{toc}{chapter}{#1}}

% =====================================================================
\begin{document}
\pagenumbering{roman}

% ---------------------------------------------------------------------
% Title page
% ---------------------------------------------------------------------
\begin{titlepage}
\centering
\singlespacing
{\bfseries HR-ETE-GNN: Robust Volatility Forecasting via\\
Hurst-Regime-Adaptive Rényi Effective Transfer Entropy\\
Graph Neural Networks\par}
\vspace{2cm}
A Thesis\\[0.5em]
Presented to the\\
Department of Information Systems\\
and Computer Science\\
Ateneo de Manila University\par
\vspace{2cm}
In Partial Fulfillment\\
of the Requirements for the Degree\\
Bachelor of Science in Computer Science\par
\vspace{2cm}
by\\[0.5em]
{\bfseries Lorenzo Uriel L. Ramoy}\par
\vspace{1cm}
{\bfseries Advisor: Dr. John Paul Vergara}\par
\vfill
2026
\end{titlepage}
\setcounter{page}{2}

% ---------------------------------------------------------------------
% Abstract
% ---------------------------------------------------------------------
\frontchapter{ABSTRACT}

Accurate volatility forecasting is critical for portfolio optimization,
risk management, and derivatives pricing. Recent work by Lee and Cho
(2025) demonstrated that Hurst Exponent Effective Transfer Entropy Graph
Neural Network (H-ETE-GNN) which constructs directed inter-asset graphs
from Shannon-based Effective Transfer Entropy (ETE) and adapts to regime
changes via the Hurst exponent outperforms traditional univariate and
correlation-based models for one-day-ahead realized volatility
prediction. However, the authors acknowledged that their model's
performance deteriorates during extreme market conditions, attributing
this to Shannon entropy's equal weighting of all regions of the return
distribution, which dilutes tail-event information that dominates crisis
periods. This paper addresses that limitation by replacing the
Shannon-based ETE with Effective Rényi Transfer Entropy (ER-TE) which is
a generalization grounded in the theoretical framework of Jizba,
Lavička, and Tabachová (2022) where a tunable order parameter $\alpha$
allows the measure to emphasize rare, extreme events ($\alpha < 1$) or
typical events ($\alpha > 1$), reducing to Shannon TE at
$\alpha \to 1$. We denote the resulting architecture
\textbf{HR-ETE-GNN} (Hurst + Rényi Effective Transfer Entropy Graph
Neural Network). Using daily price data of 10 major country-level ETFs
over 19 years (2003--2024), Forecasting performance is assessed via
QLIKE as the primary metric for comparing model performance.

% ---------------------------------------------------------------------
% Acknowledgements
% ---------------------------------------------------------------------
\frontchapter{ACKNOWLEDGEMENTS}

I would like to thank Dr. John Paul Vergara for guiding me throughout
the whole thesis process through his advice and consultations. I would
also extend my gratitude to Dr. Clark Go for helping me with this paper
by providing me with the resources and advice in understanding the
concepts.

% ---------------------------------------------------------------------
% Lists
% ---------------------------------------------------------------------
\begingroup
\singlespacing
\tableofcontents
\clearpage
\phantomsection\addcontentsline{toc}{chapter}{LIST OF FIGURES}
\listoffigures
\clearpage
\phantomsection\addcontentsline{toc}{chapter}{LIST OF TABLES}
\listoftables
\endgroup

\clearpage
\pagenumbering{arabic}

% =====================================================================
\chapter{INTRODUCTION}\label{ch:intro}
% =====================================================================

\section{Context of the Study}

Volatility prediction in financial markets is important when it comes to
risk management \cite{andersen2003,engle2004}. Volatility refers to the
magnitude of fluctuations in an asset's returns over a given time
horizon. Formally, it is most commonly operationalized as the standard
deviation of asset returns and serves as the measure of financial risk
in modern asset pricing, and portfolio construction
\cite{andersen2003,engle2004}. In recent years, different models were
employed to forecast volatility ranging from traditional models like
Generalized Autoregressive Conditional Heteroskedasticity (GARCH)
\cite{bollerslev1986} to deep learning models like Long Short-Term
Memory (LSTM) \cite{hochreiter1997}. Among the latter, Graph Neural
Networks (GNNs) have been used as a model for financial time series
because of their capabilities to model interdependencies among assets in
a graph, where nodes represent individual securities and edges encode
the strength and direction of inter-asset relationships \cite{chen2023}.

The graph tells the network which assets influence which. Granger
Causality introduces directionality but assumes linearity and
stationarity \cite{granger1969}. Transfer Entropy (TE), introduced by
Schreiber \cite{schreiber2000}, offers a model-free, non-parametric
measure of directional information flow that captures both linear and
non-linear dependencies, making it theoretically superior for financial
data characterized by fat tails, volatility clustering, and asymmetric
dependence structure \cite{behrendt2019}.

Lee and Cho \cite{lee2025} proposed a hybrid model that further advances
the proposed models with the H-ETE-GNN model which integrates three
components: (1) Effective Transfer Entropy that corrects a positive bias
that Transfer Entropy creates introduced by Marschinski and Kantz
\cite{marschinski2002}. This was used to quantify information flow among
ten country exchange-traded funds (ETFs). (2) A Graph Convolutional
Network (GCN). It is then used to construct graph convolutional networks
utilizing ETE-based graphs to forecast one-day-ahead realized
volatility. (3) A Hurst-exponent-based regime detection mechanism that
retrains the graph when the market transitions between persistent and
anti-persistent states. Their empirical results over a 19-year sample
(2003--2024) demonstrated significant improvements over LSTM, GRU, and
GNN baselines using Granger Causality, Pearson Correlation, or standard
Transfer Entropy as edge metrics.

However, Lee and Cho \cite{lee2025} identified a limitation with their
hybrid model. The model's performance degraded during extreme market
events. They attributed this degradation to two factors: the Hurst
exponent operating under a monofractal assumption, which may not fully
capture the asymmetric and complex nature of the market structure; and
the limitation of ETE in detecting extreme, localized nonlinear
interactions, potentially missing crucial regime changes. This has led
them to recommending algorithms that can detect event significance or
extreme probabilities. One of their recommendations to incorporate
extreme probabilities is Rényi-entropy-based Transfer Entropy to better
capture key information flow, increasing its robustness and adaptability
across diverse market regimes, including crises.

\section{Research Objectives}

The objective of this research is to develop and evaluate HR-ETE-GNN, a
Hurst regime-adaptive graph neural network that employs Effective Rényi
Transfer Entropy (ER-TE) in place of Shannon-based Transfer Entropy for
one-day-ahead realized-volatility (RV) forecasting across country-level
iShares MSCI ETFs, with the aim of improving forecast robustness under
extreme market regimes while preserving calm-period performance. This
study operationalizes Lee and Cho's \cite{lee2025} recommendations of
their H-ETE-GNN hybrid model and the theoretical results of Rényi
Transfer Entropy of Jizba, Lavička, and Tabachová \cite{jizba2022}.

In support of this, the specific objectives are as follows:

\begin{enumerate}
  \item Design and implement an ER-TE based adjacency construction using
        Rényi entropy with shuffled-surrogate bias correction that
        substitutes Shannon ETE in the H-ETE-GNN pipeline.
  \item To evaluate whether HR-ETE-GNN based graph construction yields
        statistically significant reduction in expected forecast loss
        relative to H-ETE-GNN.
  \item To compare ER-TE and Shannon ETE network topologies across crisis
        and non-crisis windows, identifying whether tail-sensitive
        information flow changes which countries are dominant
        transmitters of volatility information.
\end{enumerate}

\section{Research Questions}

\begin{enumerate}
  \item Does replacing Shannon-Based Effective Transfer Entropy (ETE)
        with Effective Rényi Transfer Entropy (ER-TE) as the graph
        construction metric in the Graph Neural Network of the H-ETE-GNN
        framework:
        \begin{enumerate}[label=(\alph*)]
          \item Produce a statistically significant reduction in expected
                QLIKE loss for one-day-ahead realized-volatility forecasts
                relative to H-ETE-GNN
        \end{enumerate}
  \item When the directed weighted adjacency matrix is reconstructed from
        ER-TE versus Shannon ETE on the same crisis and non-crisis
        windows;
        \begin{enumerate}[label=(\alph*)]
          \item Does the resulting topology differ in a measurable way --
                specifically, do the \textbf{top volatility information
                transmitters change} identity under tail sensitive
                $\alpha < 1$ measurement during crisis periods relative to
                Shannon measurement
        \end{enumerate}
\end{enumerate}

\section{Scope and Limitations}\label{sec:scope}

\subsubsection*{Scope}

The analysis is restricted to ten country-level iShares MSCI ETFs (EWA,
EWC, EWG, EWJ, EWT, EWU, EWW, EWY, EWZ, and EZA), observed at daily
frequency from February 2003 through May 2024. The MSCI World ETF is
used exclusively as a global regime proxy for Hurst-based detection and
is not included in the forecasting target.

In advancing the H-ETE-GNN hybrid model, only the entropy basis used for
graph construction is changed. The multi-scale 1D convolution, three GCN
layers, MSE loss, Adam optimizer, and early-stopping criteria are held
identical to Lee and Cho \cite{lee2025}. ER-TE is estimated with
$k = 1$, an $\alpha$-grid of $\{0.3, 0.5, 0.7, 0.9, 1.0, 1.5, 2.0\}$, and
$m$-shuffled surrogates for bias correction. Comparisons are drawn
against ETE-GNN, TE-GNN, Granger-GNN, Pearson-GNN, GRU, LSTM, and
H-ETE-GNN as the primary baseline.

The study does not consider intraday data, alternative
realized-volatility measures, sector or single-stock universes,
attention-based architectures such as GAT, or adaptive/regime-contingent
$\alpha$ selection.

\subsubsection*{Limitations}

\begin{enumerate}
  \item A vanilla GCN without attention is retained to preserve
        cleanliness, and a GAT-style mechanism is deferred to avoid
        confounding attribution.
  \item The k-NN Rényi Transfer estimator exhibits higher variance at
        extreme $\alpha$ values and in higher dimensions; results at
        $\alpha = 0.3$ fall below the theoretical $\alpha$-Gaussian
        validity lower bound of $\alpha > 0.4$ for $k = 1$ and are
        reported with caution.
  \item The use of a single global regime proxy (MSCI World) presumes a
        synchronized regime shift across all countries and may fail to
        capture regional differences.
  \item Comparisons against the secondary benchmarks are descriptive. The
        results do not establish that HR-ETE-GNN is superior to any model
        other than H-ETE-GNN.
\end{enumerate}

\section{Significance of Study}

The study directly operationalizes the stated recommendations of Lee and
Cho \cite{lee2025} that non-Shannon entropies are to be explored within
the H-ETE-GNN framework by providing a calibrated $\alpha$-grid and a
regime-stratified evaluation protocol. In bridging non-Shannon
entropies, we implement the information-theoretic framework of Jizba,
Lavička, and Tabachová \cite{jizba2022} on Rényi Transfer Entropy with
the H-ETE-GNN model. By grounding graph construction in a tunable-order
entropy, the work supplies a principled rationale for why tail-sensitive
$\alpha < 1$ should matter precisely when returns deviate from
Gaussianity, and why it should be immaterial when they do not.

The study also provides a structural comparison between the topology of
the networks of a Shannon-ETE and an ER-TE network, revealing whether
different countries emerge as dominant volatility transmitters under
tail-sensitive measurement.

Lastly, the realized-volatility forecasts are direct inputs to different
finance applications like Value-at-Risk, options pricing, and dynamic
hedging. If an ER-TE-based graph construction improves crisis-period RV
forecasts, this could have direct implications for approaches in
portfolio-level risk limits.

% =====================================================================
\chapter{REVIEW OF RELATED LITERATURE}\label{ch:rrl}
% =====================================================================

\section{Theoretical Framework}

\subsection{Shannon Entropy and Transfer Entropy}

Let $X = \{x_i\}$ be a discrete random variable with probability
distribution $P(X)$. The Shannon entropy of $X$ is defined as:
\[
  H(X) = -\sum_{x} p(x)\log_{2} p(x)
\]

Shannon entropy measures the average uncertainty (in bits) associated
with the outcome of $X$. Mutual information, on the other hand,
quantifies the statistical dependence of two variables. The mutual
information for two random variables $X$ and $Y$, denoted by $I(X;Y)$,
can be seen as follows:
\[
  I(X;Y) = H(X) + H(Y) - H(X,Y) = H(X) - H(X \mid Y)
\]
Where $H(X)$ is the Shannon entropy of a random variable $X$ and
$H(X \mid Y)$ denotes the conditional entropy representing the
uncertainty of $X$ given $Y$.

The concept of Mutual Information (MI) is extended to Shannon Transfer
Entropy (TE) by conditioning on $X$'s own past \cite{schreiber2000}. TE
measures the flow of information between two time series $X$ and $Y$ by
measuring how much the past of $Y$ reduces uncertainty about the next
value of $X$, beyond what $X$'s own past already provides:
\[
  TE_{Y \rightarrow X}^{(k,l)}
    = H\bigl(X_{t+1} \mid X_{t}^{(k)}\bigr)
    - H\bigl(X_{t+1} \mid X_{t}^{(k)}, Y_{t}^{(l)}\bigr)
\]
\[
  TE_{Y \rightarrow X}^{(k,l)}
    = \sum_{i} p\bigl(x_{t+1}, x_{t}^{(k)}, y_{t}^{(l)}\bigr)
      \log\frac{p\bigl(x_{t+1} \mid x_{t}^{(k)}, y_{t}^{(l)}\bigr)}
               {p\bigl(x_{t+1} \mid x_{t}^{(k)}\bigr)}
\]
Where $X_{t}^{(k)} = (x_{t}, x_{t-1}, \ldots, x_{t-k+1})$ and
$Y_{t}^{(l)} = (y_{t}, y_{t-1}, \ldots, y_{t-l+1})$ denote the past $k$
values of $X$ and past $l$ values of $Y$. The conditional probability
$p(x_{t+1} \mid x_{t}^{(k)}, y_{t}^{(l)})$ measures how much better the
probability of $x_{t+1}$ is predicted when we incorporate the past
information of $Y$ along with $X$. The sum is then calculated over all
possible joint combinations of $(x_{t+1}, x_{t}^{(k)}, y_{t}^{(l)})$. TE
is then simplified to:
\[
  TE_{Y \rightarrow X}
    = \sum_{i} p(x_{t+1}, x_{t}, y_{t})
      \log\frac{p(x_{t+1} \mid x_{t}, y_{t})}{p(x_{t+1} \mid x_{t})}
\]

\subsection{Effective Transfer Entropy}

TE, although it captures nonlinear and asymmetric dependencies, can lead
to a positive bias where TE cannot be negative, finite-sample noise
accumulates in one direction only, leaving a strictly positive
measurement even between unrelated series when estimated from finite
samples. Marschinski and Kantz \cite{marschinski2002} introduced
Effective Transfer Entropy (ETE) to correct this:
\[
  \ETE(X \rightarrow Y)
    = TE_{\text{observed}}(X \rightarrow Y)
    - \operatorname{mean}\bigl(TE_{\text{shuffled}}(X \rightarrow Y)\bigr)
\]

Shuffling destroys the temporal structure of $Y$ while preserving its
marginal distribution. The $TE_{\text{shuffled}}(X \rightarrow Y)$
estimates the finite-sample bias and the statistical significance is
assessed via a Z-score:
\[
  Z = \frac{TE_{X \rightarrow Y} - \mu_{\text{shuffled}}}{\sigma_{\text{shuffled}}}
\]
Only ETE values with $Z > 1.96$ are retained in the adjacency matrix;
all others are set to zero, yielding a sparse, directed graph.

\subsection{Rényi Entropy}

Rényi entropy \cite{renyi1961} for order $\alpha$ with probability
distribution $P$ on a finite sample $X$ is defined as:
\[
  H_{\alpha}(P) = \frac{1}{1-\alpha}\log\sum_{x \in X} p^{\alpha}(x)
\]

Rényi entropy (RE) corresponds to the average number of bits needed to
optimally encode the discrete source $X$ with probability $P(X)$
\cite{jizba2022}. For $\alpha > 1$, RE depends more on the probabilities
of the more probable values and less on the improbable ones. Meanwhile,
for $0 < \alpha < 1$, extreme events are accentuated with decreasing
$\alpha$. By choosing different values of $\alpha$, we can focus on
specific regions of the distribution, critically on the tail events that
drive crisis-period dynamics in financial markets.

\subsection{Rényi Transfer Entropy}

Following Jizba, Lavička, and Tabachová \cite{jizba2022}, the Rényi
Transfer Entropy of order $\alpha$ from $Y$ to $X$ is defined as:
\[
  T_{\alpha, Y \rightarrow X}^{R}(k,l)
    = H_{\alpha}\bigl(X_{t_{n+1}} \mid X_{n}^{(k)}\bigr)
    - H_{\alpha}\bigl(X_{t_{n+1}} \mid X_{n}^{(k)}, Y_{n}^{(l)}\bigr)
    = I_{\alpha}\bigl(X_{t_{n+1}} : Y_{n}^{(l)} \mid X_{n}^{(k)}\bigr)
\]
Where $H_{\alpha}(X \mid Y)$ is the conditional entropy of order $\alpha$
and $I_{\alpha}(X:Y)$ is the mutual information of order $\alpha$.
Explicitly:
\[
  H_{\alpha}(X \mid Y)
    = \frac{1}{1-\alpha}\log_{2}
      \frac{\sum_{x \in X, y \in Y} p^{\alpha}(x,y)}{\sum_{y \in Y} p^{\alpha}(y)}
\]
\[
  I_{\alpha}(X:Y)
    = \frac{1}{1-\alpha}\log_{2}
      \frac{\sum_{x \in X, y \in Y} p^{\alpha}(x)\,p^{\alpha}(y)}
           {\sum_{x \in X, y \in Y} p^{\alpha}(x,y)}
\]

Its difference from Shannon TE lies in three components: (1)
\textbf{Tail sensitivity}, the escort distribution $p^{\alpha}$
reweights the joint probability inside the log, allowing $\alpha < 1$ to
amplify rare-event contributions. (2) \textbf{Can be negative}, unlike
Shannon TE, RTE can take negative values. This occurs when knowing $Y$'s
past actually increases uncertainty about $X$ under the Rényi measure.
(3) \textbf{Shannon limit}, at $\alpha = 1$, Rényi TE reduces to Shannon
TE via L'Hôpital.

\subsection{Key Theorems}\label{sec:theorems}

\textbf{Jizba et al.'s Theorem 3.1 (Gaussian Equivalence).} If the joint
process $(X_{t}, Y_{t})$ is Gaussian then there is an exact equivalence
between Granger Causality and RTE, namely:
\[
  F_{Y \rightarrow X}^{(k,l)} = 2\,T_{\alpha, Y \rightarrow X}^{R}(k,l)
\]
Where $F_{Y \rightarrow X}^{(k,l)}$ is the standard Granger Causality
measure. This result is $\alpha$-independent, meaning that for Gaussian
data, the choice of $\alpha$ is immaterial --- all values give the same
causal inference. This guarantees that the Rényi extension does not
distort inference during calm periods when returns are approximately
Gaussian \cite{jizba2022}.

\textbf{Jizba et al.'s Theorem 3.2 (Heavy-tail Correction).} If the
joint process $(X_{t}, Y_{t})$ is $\alpha$-Gaussian with
$\alpha \in \bigl(\frac{1+k+l}{3+k+l}, 1\bigr]$ (i.e., finite covariance
matrix region), then
$F_{Y \rightarrow X}^{(k,l)} = 2\,T_{\alpha, Y \rightarrow X}^{R}(k,l)$
is a monotonically decreasing function of $\alpha$ (at fixed $k$ and
$l$) with zero reached at $\alpha = 1$. The leading correction is:
\[
  F_{Y \rightarrow X}^{(k,l)}
    = 2\,T_{\alpha, Y \rightarrow X}^{R}(k,l)
    + \frac{l(\alpha - 1)^{2}}{4}
    + O\bigl((\alpha - 1)^{3}\bigr)
\]

This theorem is relevant to financial returns as they exhibit heavy
tails, meaning they deviate from Gaussianity. For such data, Shannon TE
($\alpha = 1$) and Granger Causality give the same answer, but moving
$\alpha < 1$ captures additional tail-driven information flow that the
Gaussian Equivalence misses. The correction term $l(\alpha - 1)^{2}/4$
quantifies exactly how much additional information is extracted
\cite{jizba2022}.

\subsection{Effective Rényi Transfer Entropy}

Similar to the Shannon case, we define Effective Rényi Transfer Entropy
by subtracting the RTE computed on surrogate time series:
\[
  \ERTE_{\alpha, Y \rightarrow X}
    = T_{\alpha, Y \rightarrow X}^{R}
    - T_{\alpha, Y_{\text{surrogate}} \rightarrow X}^{R}
\]
Where $Y_{\text{surrogate}}$ is a phase-randomized surrogate series that
preserves the marginal distribution and autocorrelation of $Y$ but
destroys its temporal relationship with $X$. This construction,
following Jizba et al.~\cite[Eq.~49]{jizba2022}, removes finite-sample
bias while preserving the tail-sensitivity of the Rényi measure.

The balance of effective transfer entropy further quantifies the net
direction of information flow:
\[
  \ERTE_{\alpha, Y \rightarrow X}^{\text{balance}}
    = \ERTE_{\alpha, Y \rightarrow X} - \ERTE_{\alpha, X \rightarrow Y}
\]

\subsection{Hurst Exponent and Regime Adaptation}

The Hurst Exponent $H$ is estimated via Rescaled Range (R/S) analysis on
a rolling window of length $w$ \cite{hurst1951}. For a time series
segment $T_{i}$ of length $w$, the Hurst Exponent is utilized to
identify regime changes by analyzing the structure of the time series
data. It is an effective tool to interpret the complex patterns seen in
financial time series data.

Given a time series $X = \{x_{1}, x_{2}, \ldots, x_{n}\}$, we partition
these into segments of length $w$. This results in a series of segments
given by $X_{w} = \{T_{1}, T_{2}, \ldots, T_{m}\}$. We define each
segment as $T_{i} = \{x_{(i-1)w+1}, x_{(i-1)w+2}, \ldots, x_{iw}\}$,
where the total number of segments is $m = \lfloor n/w \rfloor$. By
calculating the rescaled range (R/S) values across these scales, the
scaling relationship of the time series can be determined. For a time
series segment $T_{i}$ of length $w$:

\begin{enumerate}
  \item We first compute the segment mean:
        $\overline{T}_{i} = \frac{1}{w}\sum_{j=1}^{w} x_{(i-1)w+j}$
  \item Compute the cumulative deviation:
        $Y_{k}^{(i)} = \sum_{j=1}^{k}\bigl(x_{(i-1)w+j} - \overline{T}_{i}\bigr)$
  \item Compute the range:
        $R_{i} = \max_{1 \le k \le w} Y_{k}^{(i)} - \min_{1 \le k \le w} Y_{k}^{(i)}$
  \item Compute the standard deviation:
        $S_{i} = \sqrt{\frac{1}{w}\sum_{j=1}^{w}\bigl(x_{(i-1)w+j} - \overline{T}_{i}\bigr)^{2}}$
  \item The rescaled range R/S is given by
        $\left(\frac{R}{S}\right)_{w} = \frac{1}{m}\sum_{i=1}^{m}\frac{R_{i}}{S_{i}}$
\end{enumerate}

$H$ is then estimated by regressing $\log(R/S)_{w}$ on $\log w$ via
Ordinary Least Squares (OLS).

Regime classification follows Lee and Cho (2025): using a rolling
250-day window on the MSCI World ETF (as a global proxy), a binary
regime label is assigned based on whether the proportion of days with
$H > 0.5$ in a trailing 10-day window exceeds a sensitivity threshold
$s$. When regime transition is detected, the ER-TE adjacency matrix is
recomputed and the GNN weights are fine-tuned (not fully retrained).

\subsection{Graph Convolutional Networks}

A graph $G = (V, E)$ is defined by a set of nodes $V$ and edges $E$,
encoded in an adjacency matrix $A \in \mathbb{R}^{n \times n}$, where
$A_{ij}$ measures the relationship from node $j$ to node $i$. In this
thesis, each node is one of ten country-level MSCI ETFs, and the edges
are derived from the ER-TE adjacency constructed in
Section~\ref{sec:erte}.

Classical convolutional neural networks exploit the regular grid
structure of images or time series, where translation invariance allows
a fixed-size kernel to be slid across the input. This assumption fails
for graph-structured data, since nodes have variable neighborhood sizes
and there is no canonical ordering. Graph Convolutional Networks (GCNs)
generalize the convolution operation to the irregular domain of graphs,
allowing each node's representation to be updated as a learned function
of its own features and the features of its neighbors \cite{kipf2017}.

The spectral approach, on which Kipf and Welling \cite{kipf2017} build,
defines convolution through the graph Laplacian. For an undirected graph
with adjacency $A$ and degree matrix $D$, the normalized Laplacian is
\[
  L = I_{n} - D^{-1/2} A D^{-1/2}
\]
$L$ admits an eigendecomposition of $L = U \Lambda U^{T}$, where the
columns of $U$ form a basis of graph signals analogous to the Fourier
Basis on Euclidean domains. A spectral graph convolution between a
signal $x \in \mathbb{R}^{n}$ and filter $g_{\theta}$ then takes the
form:
\[
  g_{\theta} * x = U g_{\theta}(\Lambda) U^{T} x
\]

Direct computation of this convolution is expensive because it requires
the full eigendecomposition of $L$. Defferrard, Bresson, and
Vandergheynst \cite{defferrard2016} address this by approximating
$g_{\theta}$ with a truncated expansion in Chebyshev polynomials of
order $K$ which localizes the filter to $K$-hop neighborhoods and
removes the need for eigendecomposition.

Kipf and Welling's first-order simplification \cite{kipf2017} further
simplify this Chebyshev expansion to first order $K = 1$ and apply a
renormalization trick to stabilize training yielding the layer-wise
propagation rule:
\[
  H^{(l+1)} = \sigma\bigl(\widehat{D}^{-1/2}\,\widehat{A}\,\widehat{D}^{-1/2}\,H^{(l)} W^{(l)}\bigr)
\]
Where $\widehat{A} = A + I_{n}$ adds self-loops, $\widehat{D}$ is the
degree matrix of $\widehat{A}$, $H^{(l)}$ is the matrix of node features
at layer $l$, $W^{(l)}$ is a learnable weight matrix, and $\sigma$ is a
non-linear activation. Stacking $l$ such layers allows each node to
aggregate information from neighbors up to $l$ hops away. This
formulation is the spatial interpretation of a spectral GCN: at every
layer, each node's representation is updated as a weighted average of
its neighbors' representations, with weights determined by the
(renormalized) adjacency.

\section{Direct Predecessor -- Lee and Cho (2025): H-ETE-GNN}

Lee and Cho \cite{lee2025} introduced the Hurst-Exponent Effective
Transfer Entropy Graph Neural Network (H-ETE-GNN) for one-day-ahead
realized-volatility forecasting on ten country-level iShares MSCI ETFs
over 2003--2024. Their hybrid model consisted of three components: (1)
directed, sparse adjacency matrices built from Shannon-based Effective
Transfer Entropy (ETE) with shuffled-surrogate bias correction and a
$Z > 1.96$ significance threshold; (2) a Graph Convolutional Network
(GCN) with multi-scale 1D temporal convolutions and three
graph-convolution layers; (3) a Hurst-exponent-based regime-detection
mechanism that triggers graph reconstruction and network fine-tuning
when the market transitions between persistent ($H > 0.5$) and
anti-persistent ($H < 0.5$) states. Against LSTM \cite{hochreiter1997},
GRU \cite{cho2014}, and GNN baselines using Pearson correlation or
Granger Causality \cite{granger1969} as edge metrics, H-ETE-GNN reports
statistically significant RMSE, MAE, and QLIKE gains.

Although their hybrid model performed statistically significantly better
than the aforementioned models, they identified a limitation that
motivates this paper. They reported that H-ETE-GNN's relative advantage
diminishes during the 2008 Global Financial Crisis, the 2020 COVID-19
shock, and the 2022 Russia--Ukraine episode, and they attribute this
degradation to Shannon Entropy's uniform weighting of the probability
distribution where tail events enter the ETE calculation with the same
weight as routine fluctuations, despite dominating risk during crises.
In their discussion of future work, they recommended exploring Rényi
Transfer Entropy to capture tail-driven information flow. This paper
operationalizes that recommendation through the HR-ETE-GNN model. The
present paper inherits the 10-ETF universe, the 19-year sample period,
the Hurst-based regime trigger, the multi-scale GCN architecture, the
MSE loss, the Adam optimizer, and the early-stopping criteria without
modification; only the entropy basis of the adjacency construction is
varied, so that any observed performance change can be attributed to the
Rényi generalization rather than to architectural confounders.

\section{Rényi Transfer Entropy}

\subsection{Jizba, Kleinert, and Shefaat (2012)}

Jizba, Kleinert, and Shefaat \cite{jizba2012} applied Rényi-based
information transfer to financial time series. Building on Shannon
Transfer Entropy of Schreiber \cite{schreiber2000}, the authors
introduce a Rényi-entropy generalization indexed by an order parameter
$\alpha \geq 0$ and apply it to daily returns across major international
stock indices. They establish empirically that $\alpha$ acts as a
``lens'' on the return distribution: small $\alpha$ amplifies rare tail
events, while large $\alpha$ concentrates on the distribution mode.
Shannon TE is recovered exactly as $\alpha \to 1$. Their analysis of
information flow across the DJIA, NIKKEI, FTSE, and DAX shows that Rényi
TE resolves multiscale and multifractal information-transfer patterns
that Shannon TE cannot detect, with different $\alpha$ values revealing
qualitatively different cross-market dependence structures.

This paper supplied the precedent that RTE with $\alpha \neq 1$ yields
information beyond Shannon TE specifically on financial data,
establishing that the present study's tunable-$\alpha$ design is not an
unmotivated generalization. Second, it establishes the convention of
computing RTE over an $\alpha$-grid rather than a single value. However,
this paper stopped at descriptive analysis and did not apply RTE for
forecasting. They also did not perform any bias correction via
surrogates and did not integrate RTE into any machine learning pipeline.
HR-ETE-GNN extends their work by embedding an Effective Rényi TE into a
regime-adaptive forecasting model.

\subsection{Jizba, Lavička, and Tabachová (2022)}

Jizba, Lavička, and Tabachová's \cite{jizba2022} paper on causal
inference in time series via Rényi Transfer Entropy serves as one of the
theoretical foundations on which this paper's identification rests. They
derived Rényi Transfer Entropy (RTE) from conditional Rényi entropy and
Rényi mutual information and established two results that are directly
operationalized in HR-ETE-GNN.

\textbf{Theorem 3.1 (Gaussian Equivalence)} proves that for jointly
Gaussian processes, Rényi TE equals half the Granger Causality measure
for every value of $\alpha$. This $\alpha$-independence result is the
theoretical justification for why the Rényi extension cannot distort
inference during calm periods when financial returns approximate
Gaussianity. The choice of $\alpha$ is immaterial whenever the
underlying process is well-behaved.

\textbf{Theorem 3.2 (Heavy-tail Correction)} establishes that, for an
$\alpha$-Gaussian process with power-law tails, RTE acquires a
correction term of leading order $l(\alpha - 1)^{2}/4$ on top of the
Shannon-limit value, quantifying how much additional tail-driven
information is extracted by moving below the Shannon limit.

The two theorems jointly deliver a principled argument for why
tail-emphasizing $\alpha < 1$ should matter precisely when returns
deviate from Gaussianity (crisis) and should not matter when they do not
(calm periods). They also define the Effective Rényi TE via
phase-randomized surrogates; HR-ETE-GNN adopts this construction
verbatim for bias correction.

\section{Graph Neural Network For Financial Networking}

\subsection{Kipf and Welling (2017) -- Graph Convolutional Network}

Kipf and Welling (2017) \cite{kipf2017}, introduced the Graph
Convolutional Network (GCN) which is a first-order spectral
graph-convolution layer in which each node aggregates a linear
combination of its neighbors' features at every layer, followed by a
pointwise non-linearity. The propagation rule
$H^{(l+1)} = \sigma\bigl(\widehat{D}^{-1/2}\,\widehat{A}\,\widehat{D}^{-1/2}\,H^{(l)} W^{(l)}\bigr)$,
with $\widehat{A} = A + I_{n}$ and $\widehat{D}$ the corresponding
degree matrix, reduces graph learning to sparse matrix multiplications
tractable at scale. Although the original application is semi-supervised
node classification on citation graphs, the formulation has become the
default choice wherever practitioners need a simple, well-understood
message-passing layer over a fixed adjacency.

Lee and Cho (2025) \cite{lee2025} adopt a stack of three GCN layers for
their financial application, and HR-ETE-GNN inherits this architectural
choice unchanged in order to preserve the clean-ablation protocol. The
GCN's requirement of a pre-specified adjacency matrix is precisely what
makes the entropy-basis intervention explored in this paper meaningful:
the learnable parameters of the network operate on whatever graph is
supplied, so replacing Shannon-ETE edges with ER-TE edges isolates the
contribution of the edge-construction metric.

% =====================================================================
\chapter{METHODOLOGY}\label{ch:method}
% =====================================================================

To compare the difference between the application of Rényi Transfer
Entropy and Shannon Transfer Entropy to predict stock volatility, our
methodology follows Lee and Cho's \cite{lee2025} hybrid model.

\section*{Methodology as Replicable}
\addcontentsline{toc}{section}{Methodology as Replicable}

\section{Data}

This paper will use daily closing prices of ten MSCI country-level
iShares ETFs representing both developed and emerging markets: (1) EWA,
iShares MSCI Australia, (2) EWC, iShares MSCI Canada, (3) EWG, iShares
MSCI Germany, (4) EWJ, iShares MSCI Japan, (5) EWT, iShares MSCI Taiwan,
(6) EWU, iShares MSCI United Kingdom, (7) EWW, iShares MSCI Mexico, (8)
EWY, iShares MSCI South Korea, (9) EWZ, iShares MSCI Brazil, (10) EZA,
iShares MSCI South Africa. Additionally, the iShares MSCI World ETF is
used for the Hurst-exponent-based regime detection. The data gathered
would be from February 7, 2003, to May 29, 2024.

Given $p_{t}$ as the current closing price at time $t$ of an ETF and
$p_{t-1}$ as the closing price of the day before, we compute the return
$r_{t}$ at time $t$ as the following:
\[
  r_{t} = \ln\left(p_{t} / p_{t-1}\right)
\]

Realized volatility is computed from the past $n$ trading days of
log-returns:
\[
  RV_{t} = \sqrt{\frac{1}{n}\sum_{i=1}^{n}\left(100\, r_{t-i}\right)^{2}}
\]
Where $t$ refers to the specific date, and $n$ denotes the number of
past observations. For this study, we will compute realized volatility
using the past 20 days of log-returns.

The supervised samples are formed with a look-back window of 20 trading
days and split chronologically into training (60\%), validation (20\%),
and test (20\%) blocks. All graph estimation, hyperparameter selection,
and Rényi-order selection use the training and validation blocks only;
the ER-TE adjacency matrix is estimated on training-block returns alone,
so that no test-period observation informs the graph the model uses to
forecast it.

\section{Crisis Sub-samples}

Three crisis periods are defined for regime-stratified evaluation: (1)
the Global Financial Crisis (GFC), spanning September 2008 to March
2009; (2) COVID-19, spanning February 2020 to June 2020; (3) the
Russia--Ukraine conflict, spanning February 2022 to May 2022.

These windows are defined by reference to the documented episodes of
extreme market stress themselves, and not by the behaviour of the Hurst
Exponent.

\section{ER-TE Estimation Procedure}\label{sec:erte}

\textbf{Step 1 --- Rényi entropy estimation.} We utilize the LPS
estimator of Leonenko, Pronzato, and Savani \cite{leonenko2008}, a
k-nearest-neighbor (k-NN) estimator of Rényi entropy, for a sample
$\{x_{1}, x_{2}, \ldots, x_{N}\}$ in $d$ dimensions for order
$\alpha \neq 1$:
\[
  \widehat{H}_{\alpha}
    = \frac{1}{1-\alpha}\log_{2}\left[
        C_{k} \cdot (N-1)^{1-\alpha} \cdot B_{d}^{1-\alpha}
        \cdot \frac{1}{N}\sum_{i=1}^{N}\rho_{i}^{d(1-\alpha)}
      \right]
\]
Where $\rho_{i}$ is the distance from $x_{i}$ to its $k$-th nearest
neighbor and $B_{d}$ is the volume of the $d$-dimensional unit ball.

\textbf{Step 2 --- Conditional Rényi entropy.} The conditional Rényi
entropies $H_{\alpha}(X_{t+1} \mid X_{n}^{(k)})$ and
$H_{\alpha}(X_{t+1} \mid X_{n}^{(k)}, Y_{n}^{(l)})$ are estimated using
the same k-NN approach in the appropriate joint spaces (dimensions
$d = k + 1$ and $d = k + l + 1$, respectively).

\textbf{Step 3 --- RTE computation.} For each ordered pair $(Y, X)$ of
ETFs:
\[
  T_{\alpha, Y \rightarrow X}^{R}
    = H_{\alpha}\bigl(X_{t+1} \mid X_{n}^{(k)}\bigr)
    - H_{\alpha}\bigl(X_{t+1} \mid X_{n}^{(k)}, Y_{n}^{(l)}\bigr)
\]
with lag parameters $k = 1$, following Lee and Cho \cite{lee2025}.

\textbf{Step 4 --- Effective correction.} For each pair, generate $m$
shuffled surrogates of $Y$ and compute:
\[
  \ERTE_{\alpha, Y \rightarrow X}
    = T_{\alpha, Y \rightarrow X}^{R}
    - \frac{1}{m}\sum_{i=1}^{m} T_{\alpha, Y_{i} \rightarrow X}^{R}
\]

\textbf{Step 5 --- Significance testing.} Compute the Z-score against
the shuffled distribution. Retain only ER-TE values with $Z > 1.96$ in
the adjacency matrix; set all others to zero.

\textbf{Step 6 --- Adjacency matrix.} The resulting $n \times n$ sparse,
directed, weighted matrix $A_{\alpha}$ serves as the GCN adjacency for a
given training window and Rényi order $\alpha$. Rows of $A_{\alpha}$ are
normalised to sum to one before use.

Row normalisation is a departure from the specification of Lee and Cho
\cite{lee2025}. And it is made for experimental control rather than in
spite of it. Under their node-update rule (Section~\ref{sec:gcn}) the
neighborhood term enters in proportion to the raw magnitude of the
adjacency entries, so the influence of the graph on the forecast is a
function of the scale of the edge weights and not only of their pattern.
Rényi and Shannon transfer entropies do not share a common scale: the
order $\alpha$ enters the Rényi entropy through a factor
$1/(1-\alpha)$, so ER-TE magnitudes vary systemically with $\alpha$ even
where the underlying network is changed. On raw weights, a comparison
between $\alpha = 1$ and $\alpha = \alpha^{*}$ arms would therefore
confound a difference in edge topology, which is the object of this
study. Row Normalisation removes that confound by placing every
adjacency on a common scale, and it is applied identically to both arms
and to every secondary graph model.

\textbf{Step 7 --- Network Comparison.} For each matrix, each ETF's
out-degree is computed as the number of retained outgoing edges and its
weighted out-degree as the sum of their ER-TE values. The two networks
are then compared on three quantities: (i) the identity and ordering of
the top-$K$ transmitters, $K = 3$; (ii) the number of directed edges
common to both matrices, reported as an overlap table alongside
adjacency heatmaps; and (iii) Kendall's rank correlation $\tau$
\cite{kendall1938} between the two out-degree rankings, with a 90\%
confidence interval obtained by stationary bootstrap \cite{politis1994}
with expected block length 20 trading days. All three are reported
separately for crisis and non-crisis windows.

The confidence interval on $\tau$ is reported to convey how much of the
observed rank divergence is compatible with sampling variation. It is a
single pre-specified quantity rather than a battery of edge-level tests,
and it is presented as a descriptive aid; the language used to discuss
it in Chapters~\ref{ch:results} and~\ref{ch:conclusion} is that of
compatibility with noise, not of statistical significance.

\section{Graph Convolutional Network}\label{sec:gcn}

The GNN architecture follows Lee and Cho \cite{lee2025} and is grounded
in the spectral GCN framework of Kipf and Welling \cite{kipf2017}, to
enable a clean ablation.

\textbf{Multi-scale temporal feature extraction.} Given an input time
series of length $t$ for each of $n$ ETFs, multi-scale features are
extracted via 1D convolutional filters with kernel sizes
$k_{i} \in \{3, 5, 7\}$, each with 12 channels:
\[
  f = \bigoplus_{i=0}^{q} \ReLU\left(k_{i} \cdot v + b_{i}\right)
\]
Where $\bigoplus$ denotes concatenation across filter outputs. The result
is a feature matrix $H \in \mathbb{R}^{n \times d}$.

\textbf{Message passing.} The adjacency matrix $A$ is the ER-TE matrix.
The GCN update rule for node $i$ at layer $l + 1$ is:
\[
  h_{i}^{l+1} = \ReLU\Bigl(h_{i}^{l} W_{1}^{l}
      + \sum_{j \in N(i)} A_{ji}\, h_{j}^{l} W_{2}^{l}\Bigr),
  \qquad h^{0} = f
\]
Where $W_{1}^{l}$ and $W_{2}^{l}$ are learnable weight matrices, $A_{ji}$
is the ER-TE weight of the directed edge from node $j$ to node $i$, and
$N(i)$ is the set of neighbors of node $i$. With three GCN layers, each
node aggregates information from neighbors up to three hops away.

\textbf{Output.} The final GCN layer produces a prediction vector of
dimension $(1 \times n)$, representing the one-day-ahead
realized-volatility forecast for each of the $\mathbf{n}$ ETFs.

\section{Hyperparameter Configuration}

The Rényi order $\alpha$ is introduced as an additional hyperparameter
alongside those defined in Lee and Cho \cite{lee2025}. All
hyperparameters are tuned via Bayesian Optimization.

\begin{table}[htbp]
  \centering
  \singlespacing
  \caption{Hyperparameters}\label{tab:hyperparams}
  \small
  \begin{tabular}{lll}
    \toprule
    \textbf{Hyperparameter} & \textbf{Search Space} & \textbf{Description} \\
    \midrule
    Learning Rate               & $[1 \times 10^{-2} - 1 \times 10^{-5}]$ & Adam optimizer step size \\
    Batch Size                  & $[4, 8, 16]$                        & Mini-batch size \\
    Hidden State Size (Layer 1) & $[10 - 100]$                        & GCN hidden dimension \\
    Hidden State Size (Layer 2) & $[10 - 100]$                        & GCN hidden dimension \\
    Input Size ($M$)            & $[20, 40, 60]$                      & Look-back window in days \\
    Training Period ($N$)       & $[250, 500, 750]$                   & Training window length \\
    Regime Sensitivity ($s$)    & $[6, 7, 8]$                         & Hurst threshold for regime change \\
    Shuffle Simulations ($m$)   & $[100]$                             & Number of surrogate series for ETE \\
    Rényi Order ($\alpha$)      & $\{0.3, 0.5, 0.7, 0.9, 1.0, 1.5, 2.0\}$ & Tail-sensitivity parameter \\
    \bottomrule
  \end{tabular}
\end{table}

When $\alpha = 1.0$, HR-ETE-GNN reduces exactly to the original
H-ETE-GNN \cite{lee2025}, ensuring backward compatibility.

\section{Regime-Adaptive Retraining Protocol}

Following Lee and Cho \cite{lee2025}:

\begin{itemize}
  \item Compute the daily Hurst exponent $H(t)$ on the MSCI World ETF
        using a 250-day rolling R/S window.
  \item For each day $t$, count the number of days in
        $\{t - 9, \ldots, t\}$ where $H > 0.5$.
  \item If this count exceeds the threshold $s$, classify as
        ``persistent'' regime; otherwise, ``anti-persistent.''
  \item When the regime label changes from the previous day, trigger:
        \begin{itemize}
          \item Recomputation of the ER-TE adjacency matrix on the most
                recent $N$ days.
          \item Fine-tuning of the GNN weights.
        \end{itemize}
  \item If the regime remains unchanged, the current model and adjacency
        are used for prediction without modification.
\end{itemize}

\section{Benchmark Models}

The primary comparison is HR-ETE-GNN vs H-ETE-GNN. Secondary comparisons
against other models provide context.

\begin{table}[htbp]
  \centering
  \singlespacing
  \caption{Baseline Models and their Characteristics}\label{tab:baselines}
  \small
  \begin{tabular}{lllll}
    \toprule
    \textbf{Model} & \textbf{Edge Metric} & \textbf{GNN} & \textbf{Regime Adapt} & $\boldsymbol{\alpha}$ \\
    \midrule
    HAR-RV                       & N/A (univariate)                     & No                  & No    & N/A \\
    LSTM \cite{hochreiter1997}   & N/A (univariate)                     & No                  & Yes   & N/A \\
    GRU \cite{cho2014}           & N/A (univariate)                     & No                  & Yes   & N/A \\
    GNN-Pearson                  & Pearson correlation                  & GCN \cite{kipf2017} & No    & N/A \\
    GNN-Granger                  & Granger causality \cite{granger1969} & GCN \cite{kipf2017} & No    & N/A \\
    GNN-TE                       & Shannon TE \cite{schreiber2000}      & GCN \cite{kipf2017} & No    & N/A \\
    GNN-ETE                      & Shannon ETE \cite{marschinski2002}   & GCN \cite{kipf2017} & No    & N/A \\
    H-ETE-GNN \cite{lee2025}     & Shannon ETE                          & GCN                 & Hurst & $\alpha = 1$ \\
    HR-ETE-GNN (this work)       & Rényi ETE                            & GCN                 & Hurst & Tuned \\
    \bottomrule
  \end{tabular}
\end{table}

\section{Evaluation Metrics}

The pre-registered primary loss function is QLIKE, evaluated on the
variance scale. Given the realized volatility variance $\sigma^{2}$, and
the model's forecast variance $\widehat{\sigma}^{2}$ of ETF $i$ on day
$t$, we compute QLIKE as:
\[
  r_{t,i} = \frac{\sigma_{t,i}^{2}}{\widehat{\sigma}_{t,i}^{2}}
\]
\[
  \QLIKE_{t,i} = r_{t,i} - \log r_{t,i} - 1
\]

All tests in Section~\ref{sec:hypotheses} operate on a single daily
series. For each test day $t$, and ETF $i$, the loss differential is:
\[
  d_{t} = \frac{1}{n}\sum_{i}\Bigl[\QLIKE_{t,i}(\text{H-ETE-GNN})
          - \QLIKE_{t,i}(\text{HR-ETE-GNN})\Bigr]
\]
that is, the cross-sectional average across the $n = 10$ ETFs of the
per-asset loss difference, taken first, so that each test is applied to
$T$ daily observations rather than to a pooled panel of $n \times T$. A
positive $d_{t}$ means HR-ETE-GNN incurred the smaller loss on day $t$,
and every hypothesis below is therefore one-sided in the direction
$d > 0$. Pooling the panel would be inappropriate: the ten ETFs' losses
are strongly correlated within a date, a time-series HAC estimator does
not correct for that cross-sectional dependence, and a pooled test would
consequently understate the standard error and overstate significance.

All standard errors are heteroskedasticity- and
autocorrelation-consistent, computed with the Newey--West estimator
\cite{newey1987} under a Bartlett kernel with bandwidth
$\left[4(T/100)^{2/9}\right]$, where $T$ is the number of test days; the
bandwidth actually used is reported. HAC correction is necessary rather
than precautionary here, since volatility clusters and loss
differentials are in consequence strongly serially correlated, so that
an variance estimate would manufacture significance.

\section{Hypotheses and Test Procedure}\label{sec:hypotheses}

\subsection{Hypothesis 1a --- Average predictive accuracy}\label{sec:h1a}

\textbf{$\boldsymbol{H_0}$}: Over the test period, HR-ETE-GNN and
H-ETE-GNN have equal expected loss, and any observed difference is
sampling variation.

\textbf{$\boldsymbol{H_1}$}: HR-ETE-GNN has strictly lower expected loss
than H-ETE-GNN.

To test for this we use The Giacomini--White test of conditional
predictive ability \cite{giacomini2006} in its unconditional form,
one-sided, with the HAC standard error. The statistic is
$t = \bar{d} / \mathrm{SE}(\bar{d})$ and the one-sided p-value is
$1 - \Phi(t)$.

$H_0$ is rejected at the 5\% level for $t > 1.645$. In addition to the
p-value, the mean differential $\bar{d}$, its expression as a percentage
of the baseline's mean loss, and a one-sided 95\% lower confidence bound
are reported. A statistically significant improvement smaller than 5\%
of baseline loss is reported as statistically detectable but
economically immaterial. A negative $\bar{d}$ with a large statistic is
reported as evidence that the Shannon baseline is superior; the test is
not converted to a two-sided form after the fact.

\subsection{Research Question 2 --- Descriptive network comparison}

Research Question 2 concerns the structure of the estimated networks
rather than the accuracy of the forecasts they support, and it is
addressed descriptively. No hypothesis test is attached to it, and no
claim of statistical significance is made about any structural
difference reported under it. Its findings are secondary to those of
Section~\ref{sec:h1a}.

During crisis windows, the set of dominant volatility-information
transmitters, the top-ranked source ETFs by out-degree in the retained
adjacency matrix, is expected to differ between Shannon ($\alpha = 1$)
and tail-sensitive ($\alpha^{*}$) measurement, and to coincide outside
them. The basis for this expectation is the theory of
Section~\ref{sec:theorems} and is directional: by Theorem 3.1 the two
measures agree exactly when returns are Gaussian, which is the
calm-period case, and by Theorem 3.2 they diverge as returns become
heavy-tailed.

The three quantities of Section~\ref{sec:erte}, Step 7, top-$K$
transmitter identity, edge overlap counts with adjacency heatmaps, and
Kendall's $\tau$ between the two out-degree rankings with a
stationary-bootstrap confidence interval, computed separately for crisis
and non-crisis windows.

A formal test of the null that the two transmitter rankings are
identical would require a resampling procedure over the edge set
together with correction for the ninety ordered pairs screened
simultaneously. That machinery lies outside the scope defined in
Section~\ref{sec:scope}. The consequence is stated plainly rather than
left to the reader to infer: the structural differences reported in
Chapter~\ref{ch:results} may in part reflect estimation noise, and no
significance is claimed for them. They are offered as a characterisation
of what the two measures see, and as context for the forecasting
results, which are where the inferential weight of this study rests.

Relation to the forecasting hypotheses. The structural comparison is a
claim about measurement, not about prediction, and the two are
independent by construction. The two entropy bases may describe visibly
different networks without that difference translating into forecast
accuracy under the architecture of Section~\ref{sec:gcn}, and
conversely. The two sets of results are therefore reported without
conditioning either on the other, and in particular a divergence in
network structure is not offered as evidence for, nor its absence as
evidence against, Hypotheses 1a.

% =====================================================================
\chapter{PRELIMINARY RESULTS}\label{ch:results}
% =====================================================================

\section{Preliminary Results}

The results in this section come from a reduced-scope execution of the
pipeline described in Chapter~\ref{ch:method} and are reported to
demonstrate the evaluation format rather than to establish a finding.
They should not be read as estimates of the confirmatory result.

\begin{table}[htbp]
  \centering
  \singlespacing
  \caption{Mean QLIKE by Model}\label{tab:qlike}
  \begin{tabular}{lll}
    \toprule
    \textbf{Model} & \textbf{QLIKE} & \textbf{vs.\ baseline} \\
    \midrule
    H-ETE-GNN  & 0.01360 & -- \\
    HR-ETE-GNN & 0.01430 & $-5.15\%$ \\
    \bottomrule
  \end{tabular}
\end{table}

\begin{table}[htbp]
  \centering
  \singlespacing
  \caption{Predictive Accuracy Comparison}\label{tab:gw}
  \begin{tabular}{ll}
    \toprule
    \textbf{Quantity} & \textbf{Value} \\
    \midrule
    Test Days ($T$)                        & 314 \\
    Mean loss differential ($\bar{d}$)     & $-0.0007$ \\
    Relative Change vs Baseline            & $-5.15\%$ \\
    HAC Standard Error                     & 0.000363 \\
    GW t-statistic                         & $-1.928$ \\
    One-sided p-value                      & 0.973 \\
    Decision at 5\%                        & Failed to reject \\
    \bottomrule
  \end{tabular}
\end{table}

Over $T = 314$ test days, mean QLIKE was 0.01360 for H-ETE-GNN and
0.01430 for HR-ETE-GNN at $\alpha = 0.5$, giving a mean loss
differential of $\bar{d} = -0.000700$, or $-5.15\%$ of baseline loss.
The HAC standard error was 0.000363, yielding a Giacomini--White
statistic of $t = -1.928$ and a one-sided p-value of 0.973 against the
pre-registered alternative. The null of equal predictive accuracy is not
rejected at the 5\% level. Under the decision rule of
Section~\ref{sec:h1a} a negative differential of this magnitude is
reported as favouring the Shannon baseline rather than as an absence of
difference; consistent with that rule, the test is not restated in the
opposite direction.

\begin{table}[htbp]
  \centering
  \singlespacing
  \caption{Out-Degree by Country, per $\alpha$}\label{tab:outdegree}
  \begin{tabular}{lcc}
    \toprule
    \textbf{Country} & $\boldsymbol{\alpha = 0.5}$ \textbf{(Rényi)} & $\boldsymbol{\alpha = 1}$ \textbf{(Shannon)} \\
    \midrule
    Germany      & 5 & 6 \\
    Australia    & 2 & 0 \\
    Canada       & 3 & 1 \\
    UK           & 3 & 3 \\
    South Korea  & 2 & 1 \\
    Brazil       & 0 & 0 \\
    Japan        & 1 & 1 \\
    Taiwan       & 2 & 4 \\
    Mexico       & 1 & 1 \\
    South Africa & 1 & 3 \\
    \bottomrule
  \end{tabular}
\end{table}

Germany is the leading volatility-information transmitter under both
measures, with out-degree 6 under Shannon and 5 under Rényi. The top 3
transmitters for Rényi are Germany, Canada, and the UK. Meanwhile, for
Shannon they are Germany, Taiwan, and the UK.

% =====================================================================
\chapter{CONCLUSION}\label{ch:conclusion}
% =====================================================================

% TODO: Conclusion (empty in the source manuscript).

% =====================================================================
%  Bibliography (numbered in the same order as the original manuscript)
% =====================================================================
\begingroup
\singlespacing
\begin{thebibliography}{99}
\addcontentsline{toc}{chapter}{BIBLIOGRAPHY}

\bibitem{andersen2003}
Andersen, T.G., Bollerslev, T., Diebold, F.X., and Labys, P. 2003.
Modeling and forecasting realized volatility. \emph{Econometrica} 71, 2
(2003), 579--625.

\bibitem{engle2004}
Engle, R.F. 2004. Risk and volatility: Econometric models and financial
practice. \emph{American Economic Review} 94, 3 (2004), 405--420.

\bibitem{bollerslev1986}
Bollerslev, T. 1986. Generalized autoregressive conditional
heteroskedasticity. \emph{Journal of Econometrics} 31, 3 (1986),
307--327.

\bibitem{hochreiter1997}
Hochreiter, S., and Schmidhuber, J. 1997. Long short-term memory.
\emph{Neural Computation} 9, 8 (1997), 1735--1780.

\bibitem{chen2023}
Chen, J., and Robert, C.Y. 2023. Multivariate realized volatility
forecasting with graph neural network. \emph{Proceedings of the Third
ACM International Conference on AI in Finance} (2023), 422--430.

\bibitem{granger1969}
Granger, C.W.J. 1969. Investigating causal relations by econometric
models and cross-spectral methods. \emph{Econometrica} 37, 3 (1969),
424--438.

\bibitem{schreiber2000}
Schreiber, T. 2000. Measuring information transfer. \emph{Physical
Review Letters} 85, 2 (2000), 461--464.

\bibitem{behrendt2019}
Behrendt, S., Dimpfl, T., Peter, F.J., and Zimmermann, D.J. 2019.
RTransferEntropy --- Quantifying information flow between different time
series using effective transfer entropy. \emph{SoftwareX} 10 (2019),
100265.

\bibitem{lee2025}
Lee, S., and Cho, P. 2025. Graph-based stock volatility forecasting with
effective transfer entropy and Hurst-based regime adaptation.
\emph{Fractal and Fractional} 9, 6 (2025), 339.

\bibitem{marschinski2002}
Marschinski, R., and Kantz, H. 2002. Analysing the information flow
between financial time series. \emph{European Physical Journal B} 30, 2
(2002), 275--281.

\bibitem{jizba2022}
Jizba, P., Lavička, H., and Tabachová, Z. 2022. Causal inference in time
series in terms of Rényi transfer entropy. \emph{arXiv preprint}
arXiv:2203.11407 (2022).

\bibitem{kupiec1995}
Kupiec, P.H. 1995. Techniques for verifying the accuracy of risk
measurement models. \emph{Journal of Derivatives} 3, 2 (1995), 73--84.

\bibitem{christoffersen1998}
Christoffersen, P.F. 1998. Evaluating interval forecasts.
\emph{International Economic Review} 39, 4 (1998), 841--862.

\bibitem{renyi1961}
Rényi, A. 1961. On measures of entropy and information.
\emph{Proceedings of the Fourth Berkeley Symposium on Mathematical
Statistics and Probability} 1 (1961), 547--561.

\bibitem{hurst1951}
Hurst, H.E. 1951. Long-term storage capacity of reservoirs.
\emph{Transactions of the American Society of Civil Engineers} 116
(1951), 770--799.

\bibitem{kipf2017}
Kipf, T.N., and Welling, M. 2017. Semi-supervised classification with
graph convolutional networks. \emph{Proceedings of ICLR} (2017).

\bibitem{jizba2012}
Jizba, P., Kleinert, H., and Shefaat, M. 2012. Rényi's information
transfer between financial time series. \emph{Physica A} 391, 10 (2012),
2971--2989.

\bibitem{leonenko2008}
Leonenko, N., Pronzato, L., and Savani, V. 2008. A class of Rényi
information estimators for multidimensional densities. \emph{Annals of
Statistics} 36, 5 (2008), 2153--2182.

\bibitem{cho2014}
Cho, K., van Merriënboer, B., Gulcehre, C., Bahdanau, D., Bougares, F.,
Schwenk, H., and Bengio, Y. 2014. Learning phrase representations using
RNN encoder-decoder for statistical machine translation.
\emph{Proceedings of EMNLP} (2014), 1724--1734.

\bibitem{defferrard2016}
Defferrard, M., Bresson, X., and Vandergheynst, P. 2016. Convolutional
neural networks on graphs with fast localized spectral filtering.
\emph{Advances in Neural Information Processing Systems} 29 (2016),
3844--3852.

\bibitem{politis1994}
Politis, D.N., and Romano, J.P. 1994. The stationary bootstrap.
\emph{Journal of the American Statistical Association} 89, 428 (1994),
1303--1313.

\bibitem{kendall1938}
Kendall, M.G. 1938. A new measure of rank correlation.
\emph{Biometrika} 30, 1--2 (1938), 81--93.

\bibitem{corsi2009}
Corsi, F. 2009. A simple approximate long-memory model of realized
volatility. \emph{Journal of Financial Econometrics} 7, 2 (2009),
174--196.

\bibitem{patton2011}
Patton, A.J. 2011. Volatility forecast comparison using imperfect
volatility proxies. \emph{Journal of Econometrics} 160, 1 (2011),
246--256.

\bibitem{newey1987}
Newey, W.K., and West, K.D. 1987. A simple, positive semi-definite,
heteroskedasticity and autocorrelation consistent covariance matrix.
\emph{Econometrica} 55, 3 (1987), 703--708.

\bibitem{giacomini2006}
Giacomini, R., and White, H. 2006. Tests of conditional predictive
ability. \emph{Econometrica} 74, 6 (2006), 1545--1578.

\bibitem{diebold1995}
Diebold, F.X., and Mariano, R.S. 1995. Comparing predictive accuracy.
\emph{Journal of Business \& Economic Statistics} 13, 3 (1995),
253--263.

\bibitem{harvey1997}
Harvey, D., Leybourne, S., and Newbold, P. 1997. Testing the equality of
prediction mean squared errors. \emph{International Journal of
Forecasting} 13, 2 (1997), 281--291.

\bibitem{clark2007}
Clark, T.E., and West, K.D. 2007. Approximately normal tests for equal
predictive accuracy in nested models. \emph{Journal of Econometrics}
138, 1 (2007), 291--311.

\end{thebibliography}
\endgroup

% =====================================================================
%  Appendix
% =====================================================================
\appendix
\frontchapter{APPENDIX}

% TODO: Appendix (empty in the source manuscript).

\end{document}
